\documentclass[a4paper,12pt]{article}
\usepackage{amsmath,amssymb,amsthm}
\usepackage[margin=1in]{geometry}

\newtheorem{theorem}{Theorem}
\newtheorem{lemma}[theorem]{Lemma}

\title{Problem 353: Risky Moon}
\author{Project Euler}
\date{}

\begin{document}
\maketitle

\section{Problem Statement}

A spaceship travels on the surface of a sphere from the north pole to the south pole, hopping between lattice points. For a sphere of integer radius~$r$, the lattice points are integer triples $(a,b,c)$ with $a^2+b^2+c^2 = r^2$. The risk of traveling a great-circle arc of length~$\theta$ is $(\theta/2)^2$. For each prime $r \le 14$, find the minimum total risk from $(0,0,r)$ to $(0,0,-r)$. Compute the sum of these minima, rounded to 10~decimal places.

\section{Mathematical Foundation}

\begin{lemma}[Risk reduction by splitting]
For points $P_1, Q, P_2$ on the sphere with arc half-lengths $\alpha = \theta(P_1,Q)/2$ and $\beta = \theta(Q,P_2)/2$, we have
\[
\alpha^2 + \beta^2 \le (\alpha + \beta)^2 \le \left(\frac{\theta(P_1,P_2)}{2}\right)^2\!,
\]
so inserting intermediate points can only reduce total risk.
\end{lemma}

\begin{proof}
Since $\alpha, \beta > 0$, we have $\alpha^2 + \beta^2 < (\alpha+\beta)^2$. By the spherical triangle inequality, $\alpha + \beta \ge \theta(P_1,P_2)/2$, but the total risk $\alpha^2 + \beta^2$ is strictly less than the single-hop risk $(\alpha+\beta)^2$ whenever the triangle inequality is tight. Hence subdividing arcs reduces risk.
\end{proof}

\begin{theorem}[Optimal path via Dijkstra]
For each prime~$r$, the minimum-risk path from $(0,0,r)$ to $(0,0,-r)$ is found by Dijkstra's algorithm on the complete graph of lattice points on the sphere of radius~$r$, with edge weight
\[
\rho(P_i, P_j) = \left(\frac{1}{2}\arccos\!\left(\frac{P_i \cdot P_j}{r^2}\right)\right)^{\!2}.
\]
\end{theorem}

\begin{proof}
The risk is additive, non-negative, and defined on a finite graph. Dijkstra's algorithm is correct for shortest paths in graphs with non-negative edge weights. The graph is complete (all pairwise hops are available) and finite (finitely many lattice points on a sphere of fixed radius), so the algorithm terminates and returns the global optimum.
\end{proof}

\begin{lemma}[Lattice point enumeration]
The integer solutions to $a^2+b^2+c^2 = r^2$ for a given~$r$ can be enumerated in $O(r^2)$ time by iterating $c \in [-r, r]$ and finding all $(a,b)$ with $a^2+b^2 = r^2 - c^2$.
\end{lemma}

\begin{proof}
For each of the $2r+1$ values of~$c$, the equation $a^2+b^2 = r^2-c^2$ is solved by iterating $a \in [0, \lfloor\sqrt{r^2-c^2}\rfloor]$ and checking whether $r^2-c^2-a^2$ is a perfect square. This takes $O(r)$ per value of~$c$, for a total of $O(r^2)$.
\end{proof}

\section{Algorithm}

\begin{enumerate}
    \item For each prime $r \in \{2,3,5,7,11,13\}$:
    \begin{enumerate}
        \item Enumerate all lattice points $(a,b,c)$ with $a^2+b^2+c^2 = r^2$.
        \item Build a complete graph with risk-weighted edges.
        \item Run Dijkstra from $(0,0,r)$ to $(0,0,-r)$.
    \end{enumerate}
    \item Sum the minimum risks over all primes.
\end{enumerate}

\section{Complexity Analysis}

\begin{itemize}
    \item \textbf{Time:} $O(n^2)$ per prime, where $n = r_3(r^2)$ is the number of lattice points. For $r \le 13$, $n$ is at most a few hundred.
    \item \textbf{Space:} $O(n)$ for Dijkstra's priority queue and distance array.
\end{itemize}

\section{Answer}

\[
\boxed{1.2759860331}
\]

\end{document}
